% SPDX-License-Identifier: Apache-2.0
\section{Prüfungsbasierter Abschluss und Freistellung von Lernarbeit}
\label{sec:assessment}
\begin{principle}{Status: geprüfter Konzeptentwurf, keine aktive Laufzeitfunktion}
Mastery bleibt der einzige gespeicherte Ergebniszustand pro Lernziel.
Das Profil berechnet daraus und aus fachlich geprüften Regeln, welche
Vorbereitung für einen festen Abschlussumfang entfallen kann. Diese
Freistellung erzeugt keine weitere Mastery. Die Abschlusslogik und der
direkte Prüfungszugang sind noch nicht als Laufzeitfunktion umgesetzt.
\end{principle}

Grundlage ist die \srcref{docs/concept/skill-graph/assessment-gated-completion.md}{überarbeitete
Konzeptspezifikation zum prüfungsbasierten Abschluss}.

\emph{Freistellung} bedeutet Verzicht auf gesonderte Lernarbeit im benannten
Profil. Sie ist weder eine Graphkante noch allgemeine schulische Anerkennung.
Die Anzeige unterscheidet geschaffte Ziele, für diesen Umfang entfallende
Vorbereitung und noch erforderliche Arbeit. Geschaffte Ziele besitzen
Mastery; entfallende Vorbereitung und noch erforderliche Arbeit werden
daraus und aus den Regeln berechnet, ohne zusätzliche individuelle
Lernzielzustände zu speichern.

\subsection{Fester Zielumfang statt wechselnder Sicht}
Für eine Person werden Curriculumfassung, Profilrevision, Level,
Platzierungen und gegebenenfalls autorisierte Wahlentscheidungen aufgelöst.
Sei $T_i\subseteq A$ die \emph{vollständige} erforderliche atomare
Zielmenge dieses Levels und
\[
 \varnothing\ne Q_i\subseteq T_i
\]
die Menge \emph{aller} erforderlichen Prüfungsziele. Jedes davon benötigt
Mastery und zählt damit unmittelbar als geschafft. Alternative Formen derselben
Prüfung schaffen keine zusätzlichen Zielpflichten. Echte alternative
Abschlussziele benötigen dagegen vorab einen verbindlichen Auswahlvertrag.
Eine derzeit im Curriculum enthaltene Pflicht darf nicht allein durch
Bezeichnung als optional verschwinden.

Anzeige und Level-3-Fokus sind keine Eingaben dieser Zielauflösung.
Ausblenden reduziert keine Abschlussanforderung. Ungeklärte Auswahl oder
ungültige Referenzen führen zu einer ungültigen beziehungsweise ungelösten
Berechnung, nicht zu einem leeren, angeblich vollständigen Level.
Ohne erforderliche Prüfungsziele ist dieses spezielle Profil nicht anwendbar.

\subsection{Eine atomar aufgelöste Bedarfsrelation}
Zunächst müssen die Grundlandschaften dem normativen Gültigkeitsvertrag
entsprechen und alle einschlägigen Verweise und Kompositionsregeln aufgelöst
sein. Für die Atome des aufgelösten Graphen definieren wir
\begin{equation}
 F_{\mathrm{atom}}=\{(a,b)\in A\times A\mid b\in\ReqAtoms(a)\},
 \qquad D_{\mathrm{atom}}=F_{\mathrm{atom}}^+.
 \label{eq:atomic-demand}
\end{equation}
$F_{\mathrm{atom}}$ und $D_{\mathrm{atom}}$ heißen im englischen Quelldokument
kurz $F$ und $D$. Die Indizes vermeiden Verwechslungen mit dem
Filterprädikat und der Menge der Filterdimensionen.
\emph{Nach} der atomaren Expansion nochmals die transitive Hülle zu
bilden ist wesentlich: Ein Atom innerhalb eines benötigten Clusters kann
selbst weitere Voraussetzungen besitzen.

Der vollständige, unter Voraussetzungen abgeschlossene Kontext ist
\begin{align}
 V_H&=T_i\cup\{b\in A\mid\exists t\in T_i:(t,b)\in D_{\mathrm{atom}}\},\\
 R_H&=D_{\mathrm{atom}}\cap(V_H\times V_H),\qquad H=(V_H,R_H).
 \label{eq:context}
\end{align}
Alle reflexiven Paare in $R_H$ führen zur Ablehnung dieser Berechnung.
Die zusätzliche Prüfung umfasst auch unterstützende und
\code{prerequisiteOnly}-Ziele. Ein atomarer Expansionszyklus in einem
unabhängigen, vom Level nicht erreichbaren Teil des ansonsten grundgültigen
Graphen betrifft dagegen nicht diese lokale Berechnung. Die globale
Grundgültigkeitsanforderung bleibt bestehen.

$R_H$ ist nach erfolgreicher Prüfung endlich, irreflexiv und transitiv.
Ein Pfeil zeigt weiterhin vom benötigenden Ziel auf seine Voraussetzung.
Bei $C\to B\to A$ enthält die Relation auch $(C,A)$.
Ein bereits abgeschlossenes $B$ beseitigt deshalb den noch offenen Bedarf
von $C$ an $A$ nicht.

Unterstützende Ziele gehören zum Kontext, aber nicht automatisch zu $T_i$.
Ein fremder späterer Verbraucher wird nicht aufgenommen, bloß weil er
dieselbe Grundlage benötigt. Die Kontextidentität bindet Person, Scope,
Rollen, Revisionen, Auswahl und Abhängigkeitsauflösung. Für die Berechnung
wird ein konsistenter aktueller Mastery-Stand gelesen. Daraus entsteht
keine zusätzliche gespeicherte Belegsammlung.

\subsection{Mastery und geprüfte Freistellungsregeln}
Sei $M_A$ die Menge der nach dem bestehenden Zieltypvertrag geschafften
atomaren Ziele. Die einzige individuelle Ergebniseingabe ist
\begin{equation}
 E=M_A\cap V_H.
 \label{eq:mastery-input}
\end{equation}
Eine vorhandene Prüfungs-Mastery zählt ohne weitere profilabhängige
Anerkennung. Ein angenommener Eintrittsstand, das Öffnen einer Aufgabe
oder eine fehlgeschlagene Prüfung erzeugt keine Mastery.
Die bestehenden Verträge von Orientierung, gewöhnlichen Zielen,
Gedächtniszielen und Prüfungen bleiben maßgeblich.

Die geschützte Menge $P\subseteq V_H$ enthält sämtliche Prüfungs-,
Orientierungs- und Gedächtnisziele sowie weitere nach Fach- oder
Typvertrag nicht freistellbare Ziele. Schutz erzeugt keine neue Zielpflicht;
ein optionales Prüfungsangebot wird dadurch nicht zur Pflichtprüfung.

Eine endliche Familie fachlich geprüfter Regeln $(S_r,W_r)$ erfüllt
\begin{align}
 &\varnothing\ne S_r\subseteq Q_i,
 &&W_r\subseteq V_H\setminus P,\\
 &\forall v\in W_r:\quad
   \exists q\in S_r:(q,v)\in R_H.
\end{align}
$S_r$ sind die gemeinsam zu schaffenden Prüfungsziele;
$W_r$ benennt die gewöhnlichen Lernziele, deren separate Arbeit diese
Prüfungen ersetzen dürfen. Im einfachsten Profil gibt es eine Regel mit
$S_r=Q_i$. Mehrere Regeln sind nur sinnvoll, wenn Teilgruppen der
Prüfungen die Freistellungen bereits eigenständig fachlich rechtfertigen.

Die erlaubte Menge ist
\begin{equation}
 W(E)=\bigcup_{r:\ S_r\subseteq E}W_r.
 \label{eq:enabled}
\end{equation}
Nur Mastery der erforderlichen Prüfungsziele aktiviert Regeln. Eine abgeleitete Freistellung
wird niemals so behandelt, als wäre damit eine Prüfung bestanden.
Erforderliche Ziele ohne geprüfte Freistellungsregel benötigen ihre eigene Mastery.

Eine Voraussetzungskante beweist keine Prüfungsabdeckung.
Auch \code{examData.coveredGoalIds} allein genügt nicht. Ein Review muss
Aufgaben, Rubriken, Bedingungen und die inhaltliche Begründung für jedes
$W_r$ anhand des bestehenden Erfolgsvertrags prüfen. Kann eine bestandene
Prüfung einen kritischen fehlenden Teil verdecken, bleibt das betroffene
Ziel außerhalb von $W_r$. Eine schriftliche Leistung ersetzt nicht ohne
Weiteres praktisches Handeln oder verifizierten Abruf. Das Profil führt
dafür keine zweite, strengere individuelle Erfolgsprüfung ein.

\subsection{Prüfungserfolg im bestehenden Mastery-Vertrag}
Für Prüfungsziele gilt derselbe Ergebnisgrundsatz wie für andere Lernziele:
geschafft oder noch nicht geschafft. Nach bestandener Bewertung wird
Mastery gespeichert; erst die bestätigte Speicherung trägt die
Abschlussberechnung. Ein Fehlversuch lässt vorhandene Mastery unverändert
und speichert weder einen weiteren Zielstatus noch einen erfolglosen
Punktestand. Auch bereits vorhandene Prüfungs-Mastery zählt unmittelbar.

Der bestehende Prüfungsablauf erlaubt Wiederholungen derselben Aufgabe
und bespricht nach der Bewertung Ergebnis und Lösung. Eine bekannte
Lösung kann die Aussagekraft eines späteren Versuchs verringern. Das ist
eine Grenze dieses Ablaufs, keine zusätzliche individuelle
Anerkennungsbedingung. Das Abschlussprofil führt weder eine zweite
Prüfungsqualifikation noch eine gespeicherte Versuchs-, Formen- oder
Lösungsexpositionshistorie ein. Es behauptet keine beaufsichtigte oder
identitätsverifizierte Zertifizierung.

Die fachliche Eignung von Aufgaben, Bewertungsrastern und Freistellungen
wird beim Entwurf und Review der statischen Regeln geprüft. Eine für
dieses Profil zu schwache Prüfung verlangt eine bessere Aufgaben- oder
Regeldefinition; sie rechtfertigt keine nachträgliche zweite Erfolgshürde
für ein bereits geschafftes Lernziel.

\subsection{Freistellung als kleinster Fixpunkt}
Alle direkten und indirekten Verbraucher eines Ziels $v$ sind
\[
 \Consumers_H(v)=\{u\in V_H\mid(u,v)\in R_H\}.
\]
Ausgehend von $E$ wird freistellbare Arbeit nur dann aufgenommen,
wenn alle ihre Verbraucher bereits durch Mastery abgeschlossen oder
selbst legitim freigestellt sind:
\begin{align}
 X_0&=E,\\
 X_{n+1}&=X_n\cup
  \{v\in W(E)\mid \Consumers_H(v)\subseteq X_n\},\label{eq:waiver-step}\\
 X_*&=\bigcup_{n\geq0}X_n,\qquad I=X_*\setminus E.
 \label{eq:waiver-result}
\end{align}
$I$ ist ausschließlich die berechnete kontextuelle Freistellungsmenge;
auch $X_*$ ist kein gespeicherter Lernzielzustand.

\paragraph{Existenz und kleinster Fixpunkt.}
Der Operator
$\Phi_E(X)=X\cup\{v\in W(E)\mid \Consumers_H(v)\subseteq X\}$
ist monoton und inflationär. Auf der endlichen Potenzmenge von $V_H$
stabilisiert die Folge nach höchstens $|V_H\setminus E|$ tatsächlich
erweiternden Runden. Ist $Y$ irgendein Fixpunkt mit $E\subseteq Y$,
folgt induktiv $X_n\subseteq Y$ für alle $n$. Also ist $X_*$ der
kleinste Fixpunkt, der $E$ enthält.

\paragraph{Schutz und offene Verbraucher.}
Jedes neu aufgenommene Element liegt in $W(E)\subseteq V_H\setminus P$,
also $I\cap P=\varnothing$. Für jedes $v\in I$ waren alle Verbraucher
schon bei seiner Aufnahme abgeschlossen oder freigestellt. Folglich
\begin{equation}
 v\in I\Longrightarrow \Consumers_H(v)\subseteq X_*.
\end{equation}
Damit kann kein freigestelltes Ziel Voraussetzung eines noch offenen
Kontextziels sein. Die Bindung jeder Regel an eine Prüfung verhindert
eine Freistellung isolierter Ziele durch eine leere Verbrauchermenge.

\paragraph{Mehr Mastery bei unverändertem Vertrag.}
Aus $E\subseteq E'$ folgt zunächst $W(E)\subseteq W(E')$ und durch
Induktion $X_*(E)\subseteq X_*(E')$. Die Menge $I$ selbst muss nicht
wachsen: Ein zuvor freigestelltes Ziel kann eigene Mastery erhalten.
Bei verändertem Scope, geänderten Regeln oder einer Mastery-Korrektur
mit zurückgenommenem Erfolg gilt diese Monotoniebehauptung nicht.

Die Rechnung kann äquivalent in topologischer Reihenfolge vom Verbraucher
zur Voraussetzung erfolgen. Für die Nachvollziehbarkeit nennt jede
Freistellung eine aktivierte Regel, die geschafften Prüfungsziele und die
bereits aufgelösten Verbraucher. Diese Erklärung wird aus denselben
Eingaben berechnet. Rechenrunden sind keine Lernstunden.

Künftig berechnet das Backend Umfang, Mastery, Abhängigkeiten und Fixpunkt
gemeinsam für alle Oberflächen und Coaches. Berechnungsgültigkeit,
Abschlussstatus und Verfügbarkeit der noch offenen Prüfungen werden
getrennt ausgewiesen. Eine fehlende freigegebene Aufgabe ist eine
Versorgungslücke und kein Lernendenversagen.

\subsection{Bedarf ist nicht dasselbe wie fehlende Beherrschung}
Abgeleitete Bedarfszustände lauten
\begin{equation}
 \begin{split}
 \Discharged_H&=\{(u,v)\in R_H\mid u\in X_*\},\\
 \Active_H&=R_H\setminus \Discharged_H.
 \end{split}
\end{equation}
Ein offenes $C$ hat weiterhin aktiven Bedarf an $A$, auch wenn $A$
ausdrücklich beherrscht ist und die gewöhnliche Voraussetzung bereits
erfüllt. Umgekehrt impliziert $A\in I$ weder $A\in M_A$ noch
$\Sat(A,M_A)$. Es werden keine Kanten gelöscht und keine Ergebnisse
geschrieben. Der Overlaystatus und die normale Frontier sind verschiedene
Auswertungen derselben fachlichen Ausgangslage.

Freistellungen erzeugen keine Lernerfolgsereignisse, Wochenquotenpunkte
oder nachgewiesenen Lernstunden und ändern keine historische Planbaseline.
Die Reduktion geplanter Restarbeit braucht einen ausdrücklichen Vertrag.
Struktur- und Aufwandstests bleiben erforderlich, ebenso geeignete
Aufgabenreviews, zuverlässige Bewertung und Lernendenerprobung.

\subsection{Durchgerechnetes Beispiel mit geteilter Grundlage}
Die erforderlichen, geschützten Prüfungen B und C benötigen A;
A benötigt D. Es gilt
\begin{align*}
 T_i=V_H&=\{A,B,C,D\},\qquad Q_i=P=\{B,C\},\\
 R_H&=\{(B,A),(C,A),(A,D),(B,D),(C,D)\},\\
 (S_r,W_r)&=(\{B,C\},\{A,D\}).
\end{align*}

\begin{figure}[htbp]
\centering
\begin{tikzpicture}[node distance=11mm]
\begin{scope}
\node[info] at (2.3,2.3) {Nur B hat Mastery};
\node[goal,fill=green!9] (b1) at (0,1.3) {B: Mastery};
\node[goal,fill=orange!12] (c1) at (0,-.3) {C: offen};
\node[goal,fill=orange!12] (a1) at (2.5,.5) {A: offen};
\node[goal,fill=orange!12] (d1) at (4.6,.5) {D: offen};
\draw[req] (b1)--(a1);
\draw[req,draw=warm,dashed] (c1)--(a1);
\draw[req,draw=warm,dashed] (a1)--(d1);
\node[info,text width=60mm] at (2.3,-1.25)
 {$W(E)=\varnothing$; keine Freistellung.};
\end{scope}
\begin{scope}[yshift=-5cm]
\node[info] at (2.3,2.3) {B und C haben Mastery};
\node[goal,fill=green!9] (b2) at (0,1.3) {B: Mastery};
\node[goal,fill=green!9] (c2) at (0,-.3) {C: Mastery};
\node[goal] (a2) at (2.5,.5) {A: frei\\Runde 1};
\node[goal] (d2) at (4.6,.5) {D: frei\\Runde 2};
\draw[req] (b2)--(a2);
\draw[req] (c2)--(a2);
\draw[req] (a2)--(d2);
\node[info,text width=95mm] at (2.3,-1.25)
 {$W(E)=\{A,D\}$; die Rechnung berücksichtigt\\auch $B\to D$ und $C\to D$.};
\end{scope}
\end{tikzpicture}
\caption[Freistellung aus Prüfungs-Mastery]{Abgeleiteter Bedarf, Pfeile vom Verbraucher zur Voraussetzung.
Gezeigt sind nur unmittelbare Beziehungen. Durchgezogen: im Ergebnis
entlasteter Bedarf; gestrichelt: aktiver Bedarf. Freistellung erzeugt keine Mastery.}
\label{fig:waivers}
\end{figure}
\FloatBarrier

Tabelle~\ref{tab:waivers} zeigt die Ergebnisse für verschiedene Mastery-Mengen.
\begin{table}[htbp]
\centering
\begin{tabularx}{\linewidth}{@{}lllX@{}}
\toprule
$E$ & $W(E)$ & $I$ & Ergebnis\\
\midrule
$\varnothing$ & $\varnothing$ & $\varnothing$ & Alles offen.\\
$\{B\}$ & $\varnothing$ & $\varnothing$ & C fehlt; keine aktivierte Regel.\\
$\{B,C\}$ & $\{A,D\}$ & $\{A,D\}$ & Erst A, anschließend D freigestellt.\\
$\{A\}$ & $\varnothing$ & $\varnothing$ & Gewöhnlicher Erfolg aktiviert keine Regel.\\
$\{A,B,C\}$ & $\{A,D\}$ & $\{D\}$ & A behält seine Mastery.\\
\bottomrule
\end{tabularx}
\caption{Vorhandene Mastery und berechnete Freistellungen.}
\label{tab:waivers}
\end{table}

Braucht ein weiteres offenes Ziel Z desselben Kontextes D, bleibt D offen.
Gehört Z ausschließlich einem späteren Level an, blockiert es die hiesige
Freistellung nicht. Ist D ein geschütztes Gedächtnisziel, darf es überhaupt
nicht in $W_r$ stehen.

\paragraph{Unterstützende Ziele können trotzdem blockieren.}
Seien $T_i=\{Q,A\}$, $Q_i=\{Q\}$ und $Q\to B\to A$.
Mit $E=\{Q\}$ und $(S_r,W_r)=(\{Q\},\{A\})$ ist B zwar kein
Ziel im Abschlussnenner, aber offener Verbraucher von A.
Deshalb gilt $I=\varnothing$. Mastery für A beendet
das Level; alternativ ermöglicht Mastery für B die Freistellung von A.
Die bloße Unterstützungsrolle berechtigt nicht zur Freistellung von B.

Im Gegenfall $Q\to A\to B$ mit denselben Zielmengen, Mastery und derselben
Regel ist B kein Verbraucher von A. A wird freigestellt und das Level ist
abgeschlossen, obwohl das unterstützende B offen bleibt; auch Schutz von B
ändert daran nichts. Würde B zusätzlich zum erforderlichen Ziel, bliebe der
Abschluss dagegen offen. Maßgeblich ist die gerichtete Bedarfsrelation:
Ein im Fixpunkt offener Verbraucher B mit $(B,A)\in R_H$ verhindert die Freistellung
eines sonst berechtigten A. Unterstützungsstatus allein blockiert keinen Abschluss.

\paragraph{Transitive Verbraucher dürfen nicht abgeschnitten werden.}
Wenn C offen ist, C B benötigt und B A benötigt, dann bleibt A wegen C
offen, auch wenn B bereits Mastery besitzt und A durch eine andere
Prüfregel grundsätzlich freistellbar ist. Eine Rechnung nur mit unmittelbaren
Verbrauchern würde hier falsch entscheiden.

\subsection{Abschluss, Zugang und weitere Lernarbeit}
Das Profil ist abgeschlossen genau dann, wenn
\begin{equation}
 \Complete(i,H,E)\iff Q_i\subseteq E\ \land\ T_i\subseteq X_*.
 \label{eq:complete}
\end{equation}
Der Abschlussnenner enthält nur Ziele. Die Prüfungsklausel folgt auch aus
$Q_i\subseteq T_i\cap P$, nennt die Pflicht aber ausdrücklich.
\glqq Allein die Prüfungen genügen\grqq{} gilt nur, wenn $E=Q_i$ tatsächlich
$T_i\subseteq X_*$ ergibt. Geschützte Orientierungs- oder Gedächtnisziele
können zusätzliche Pflichten sein. Der Rechentest belegt keine fachliche
Eignung der Aufgaben.

Der direkte Zugang erlaubt bewusst gewählte Versuche für jedes
$q\in Q_i\setminus E$, auch außerhalb der normalen Frontier.
Eine bereits geschaffte Prüfung benötigt keinen weiteren Erfolgsnachweis.
Nur die didaktische Voraussetzungssperre entfällt;
Berechtigung, Scope, geeignete Aufgabenform und vorgeschriebene Bedingungen
wie Aufsicht und praktische Sicherheit gelten weiter. Versuch und Fehlschlag erzeugen keine Freistellung.
Erfolg setzt bestätigte Mastery-Speicherung voraus; die Wiederholung
eines Schreibaufrufs erzeugt keinen weiteren Lernerfolg.

Offene Arbeit ist $T_i\setminus X_*$. Reguläre Empfehlungen kommen aus
ihrer Schnittmenge mit der geltenden Frontier; direkte Prüfungen werden
gesondert angeboten. Freigestellte Arbeit bleibt freiwillig wiederholbar,
wird in diesem Profil aber nicht erneut als Pflicht empfohlen. Eine leere
Frontier rechtfertigt weder Abschluss noch automatische Fokusausweitung.
Der berechnete Levelabschluss wird getrennt vom gewöhnlichen Fokus- und
Curriculumfortschritt angezeigt. Er setzt deren Abschlusskennzeichen
nicht auf erfüllt, solange ihre eigenen Mastery-Bedingungen offen sind.

Späteres Lernen folgt Mastery und normalen Scope-Regeln. Wird eine in einem
anderen Profil entfallende Vorbereitung hier benötigt, erklärt der Coach
die unterschiedlichen Anforderungen und bietet gezielte Demonstration oder
Lernhilfe an. Einen früheren Verlauf benennt er nur, wenn dieser bekannt ist.
Freistellung erzeugt weder globale Mastery noch eine neue Sperre im Normalmodus.
Über Beenden oder Fortsetzen entscheidet die lernende Person.

\subsection{Revisionen, Korrekturen und Einführung}
Das Ergebnis gilt für den ausgewiesenen Umfang, seine Revisionen und den
gelesenen Mastery-Stand. Ändern sich diese Eingaben, wird neu gerechnet.
Ein Verbraucher außerhalb dieses Umfangs beeinflusst die Berechnung nicht.
Eine fachlich erforderliche Ergebniskorrektur erfolgt am bestehenden
Mastery-Zustand; es gibt keinen unabhängig davon zurückziehbaren Prüfbeleg.
Ein gewöhnlicher Fehlversuch nimmt vorhandene Mastery nicht zurück.

Aus der aktuellen Mastery-Menge allein lassen sich frühere Abschlussstände
nicht zuverlässig rekonstruieren. Dieses Profil verspricht deshalb keine
unveränderliche historische Abschlussurkunde und führt dafür keine neue
Beleg-, Snapshot- oder Abschlussdatenhaltung ein. Es liefert eine aktuelle
berechnete Aussage über den benannten Lernumfang.
